\documentclass[11pt]{article}
\usepackage[a4paper,margin=2.3cm]{geometry}
\usepackage{lmodern}
\usepackage{enumitem}
\usepackage{booktabs}
\usepackage{xcolor}
\usepackage{hyperref}
\usepackage{parskip}

\definecolor{done}{HTML}{2E7D32}
\definecolor{todo}{HTML}{C62828}
\definecolor{partial}{HTML}{E65100}
\definecolor{mech}{HTML}{1565C0}

\newcommand{\done}{\textcolor{done}{\textbf{FAIT}}}
\newcommand{\todo}{\textcolor{todo}{\textbf{A FAIRE}}}
\newcommand{\parti}{\textcolor{partial}{\textbf{PARTIEL}}}
\newcommand{\mech}{\textcolor{mech}{\textbf{MECANIQUE}}}

\title{Plan de route avant soumission JASA\\[6pt]
\large\textit{Permutations Accelerate Approximate Bayesian Computation}}
\author{Memo interne --- 21 avril 2026}
\date{}

\begin{document}
\maketitle

\tableofcontents
\newpage

%% ============================================================
\section{Etat des lieux}
%% ============================================================

Le papier est a 35 pages (limite JASA), format JASA (12pt, double-spaced, \texttt{agsm.bst}, toggle \verb|\anon|). L'appendice est dans \texttt{supplement.tex} (19 pages). Les changements depuis la version Biometrika sont marques avec \verb|\old{}|/\verb|\new{}| pour relecture par Robin et Xian.

\subsection{Ce qui est fait}

\begin{itemize}[nosep]
    \item JASA formatting complet (marges, police, double-spacing, bibliographie)
    \item 35 pages atteintes (reductions : condensation OS, UM, kernels, SIR, ABC-Gibbs, Section 1 ; Fig 7+8 fusionnees ; Algorithm en \verb|\small|)
    \item \texttt{supplement.tex} separe avec numerotation S.1--S.8
    \item SW$_2$ comme metrique de qualite posterieure (remplace $\varepsilon$-vs-cost) --- repond au \textbf{Mock M1}
    \item ABC-SMC standard (sans permutations) ajoute dans les benchmarks --- repond au \textbf{Mock M2}
    \item $K\to\infty$ asymptotics integres (Proposition 1.5) --- repond partiellement au \textbf{Mock M4}
    \item $M\to\infty$ derivation deplacee en appendice S.4.1 --- repond partiellement au \textbf{Mock M4}
    \item SIR stochasticite : formule log-normal $\sigma=0.05$ documentee
    \item WABC discussion textuelle (Section 1.2, lien Wasserstein)
    \item Validation SIR synthetique : script + SBATCH + figure + texte appendice --- \textbf{code pret, pas encore execute}
    \item \texttt{jasa\_guidelines.md} cree avec toutes les consignes JASA
\end{itemize}

%% ============================================================
\section{Priorite 1 --- Bloquant pour la soumission}
%% ============================================================

\subsection{Lancer le job SIR synthetique (5 min)}
\label{sec:sir_synth}

\textbf{Statut :} \mech{} --- le code est pret, il suffit de push + submit.

\textbf{Quoi :} Le script \texttt{run\_sir\_synthetic\_validation.py} genere des donnees SIR avec parametres connus ($K=10$, 120 jours), lance permABC-SMC pour $\sigma \in \{0.01, 0.05, 0.10\}$, et sauve les resultats. Le script de figure \texttt{fig\_sir\_synthetic\_validation.py} produit le PDF pour l'appendice (Section S.8.2).

\textbf{Commandes :}
\begin{verbatim}
git add -A && git commit -m "Add SIR synthetic validation"
git push origin main
ssh ceremade-cluster "cd ~/permABC && git pull origin main"
ssh ceremade-cluster "cd ~/permABC && sbatch \
    experiments/scripts/cluster/run_sir_synthetic.SBATCH"
\end{verbatim}

\textbf{Apres execution :} recuperer les pickles, lancer le script figure, copier le PDF dans \texttt{paper/fig/}.

\textbf{Pourquoi c'est bloquant :} C'est le reproche le plus grave des reviewers (Referee 1 Biometrika, Mock M5). Sans validation sur donnees synthetiques, impossible de distinguer si permABC produit une bonne approximation d'un mauvais modele ou une mauvaise approximation.

\subsection{Hardcoder les references vers le supplement (15 min)}
\label{sec:refs}

\textbf{Statut :} \mech{} --- remplacement mecanique.

\textbf{Quoi :} Le papier principal contient 14 \verb|\ref{appendix_...}| qui pointent vers des labels definis dans \texttt{supplement.tex}. Comme les deux documents sont compiles separement, LaTeX affiche ``\textbf{??}'' dans le PDF. Il faut remplacer chaque reference par le texte en dur.

\textbf{Table de correspondance :}

\begin{center}
\begin{tabular}{ll}
\toprule
\textbf{Label} & \textbf{Remplacement} \\
\midrule
\verb|Appendix~\ref{appendix_proofs}| & Supplement, Section~S.3 \\
\verb|Appendix~\ref{sec:appendix_strat}| & Supplement, Section~S.2 \\
\verb|Appendix~\ref{appendix_os}| & Supplement, Section~S.4 \\
\verb|Appendix~\ref{appendix_os_asymptotics}| & Supplement, Section~S.4.1 \\
\verb|Appendix~\ref{appendix_um}| & Supplement, Section~S.5 \\
\verb|Appendix~\ref{appendix_kernel}| & Supplement, Section~S.6 \\
\verb|Appendix~\ref{sec:appendix_comparison}| & Supplement, Section~S.7 \\
\verb|Appendix~\ref{appendix_sir_validation}| & Supplement, Section~S.8.1 \\
\verb|Appendix~\ref{appendix_sir_synthetic}| & Supplement, Section~S.8.2 \\
\verb|Appendix~\ref{appendix_sir}| & Supplement, Section~S.8 \\
\bottomrule
\end{tabular}
\end{center}

\textbf{Attention :} si on ajoute/supprime des sections dans le supplement, les numeros changent. Faire ce remplacement en dernier, juste avant soumission.

\subsection{Theorem 1 : attenuer et etendre aux ties (1--2h, theorie)}
\label{sec:thm1}

\textbf{Statut :} \todo{} --- necessite reflexion mathematique.

\textbf{Sous-probleme A : attenuer les claims.} La phrase actuelle (Section 1.4) dit :
\begin{quote}
``the efficiency gains provided by permABC come \emph{at no cost to theoretical consistency}''
\end{quote}
C'est vrai uniquement quand $\varepsilon \to 0$. Pour tout $\varepsilon$ fini utilise en pratique, les deux pseudo-posteriors (permABC vs ABC standard) peuvent differer, sans borne connue sur l'ecart. Le mock review M3 demande explicitement : soit fournir une borne en TV ou $W_2$ pour $\varepsilon$ fini, soit attenuer.

\textbf{Proposition :} Remplacer par quelque chose comme :
\begin{quote}
``In the small-tolerance regime $\varepsilon < \varepsilon^*$, permABC and standard ABC produce identical posteriors (Theorem~1). For larger $\varepsilon$, the two pseudo-posteriors may differ; bounding this discrepancy at finite $\varepsilon$ remains an open problem.''
\end{quote}

\textbf{Sous-probleme B : etendre au cas des ties.} L'editeur de Biometrika a explicitement demande : que se passe-t-il si deux compartiments ont des donnees tres proches ? Dans ce cas $\varepsilon^*$ est tres petit et le regime $\varepsilon < \varepsilon^*$ est inaccessible.

L'idee suggeree par l'editeur : redefinir $\varepsilon^*$ comme la moitie de la distance minimale entre observations \emph{distinctes}. Alors le cardinal des permutations acceptees est constant (egal au nombre de permutations qui echangent les compartiments identiques), et Theorem~1 tient toujours avec un reweighting. La preuve est une extension mineure de Lemma~1.

\textbf{Qui :} Antoine + Robin/Xian (la preuve est mathematique).

\subsection{Formaliser les resultats $M\to\infty$ comme Propositions (1h)}
\label{sec:minfty}

\textbf{Statut :} \parti{} --- le contenu mathematique est en appendice S.4.1, mais sous forme d'equations numerotees sans enonce formel.

\textbf{Quoi :} Le mock review M4 demande des Propositions avec hypotheses explicites. Concretement :
\begin{enumerate}[nosep]
    \item Enoncer les ``mild regularity assumptions on $d$'' (continuite, compacite du support, etc.)
    \item Envelopper la convergence $d_M \to \delta_\beta$ a.s. dans une \textbf{Proposition} numerotee
    \item Envelopper le resultat \eqref{eq:os\_pushforward\_limit} (dans le main text) dans une \textbf{Proposition} numerotee, avec renvoi a l'appendice pour la preuve
    \item Ajouter un court paragraphe de preuve (ou ``proof sketch'') en appendice
\end{enumerate}

\textbf{Alternative :} Si on ne veut pas formaliser completement, etiqueter clairement comme ``heuristic'' ou ``conjectural'' avec un pointer vers un traitement rigoureux futur. Le mock review accepte cette option.

\subsection{Benchmark WABC / Hilbert / Sinkhorn vs LSA (code + cluster)}
\label{sec:wabc}

\textbf{Statut :} \todo{} --- necessite code + execution.

\textbf{Quoi :} Sur le Gaussian toy model ($K=20$, posterieure en forme close), comparer :
\begin{itemize}[nosep]
    \item permABC avec LSA exact ($\mathcal{O}(K^3)$)
    \item permABC avec Hilbert sorting ($\mathcal{O}(K\log K)$)
    \item permABC avec Sinkhorn ($\mathcal{O}(K^2 \times \text{iters})$)
    \item permABC avec warm-started swaps ($\mathcal{O}(K^2)$ amorti)
\end{itemize}

Metriques : SW$_2$ a la posterieure de reference, runtime, fraction d'iterations ou la permutation retournee coincide avec l'optimum exact.

\textbf{Pourquoi :} Mock M6 + Referee 3 (points 1 et 8). Le code des solvers existe deja dans \texttt{permabc/assignment/solvers/}. Il faut creer un script de benchmark et un SBATCH.

\textbf{Ou dans le papier :} Nouvelle figure en appendice (Section S.7 ou nouvelle section), avec une phrase dans le main text qui y renvoie.

\subsection{Resoudre les \texttt{\textbackslash old\{\}}/\texttt{\textbackslash new\{\}} (30 min, apres review Robin/Xian)}
\label{sec:oldnew}

\textbf{Statut :} \todo{} --- intentionnellement repousse.

\textbf{Quoi :} Envoyer le PDF actuel a Robin et Xian pour qu'ils valident les changements (bleu = nouveau, rouge barre = ancien). Une fois valide :
\begin{enumerate}[nosep]
    \item Accepter chaque \verb|\new{...}| : retirer le markup, garder le texte
    \item Supprimer chaque \verb|\old{...}| : retirer le bloc entier
    \item Retirer le \verb|\renewcommand{\old}[1]{}| dans \texttt{resubmission.tex}
    \item Verifier que le page count reste $\leq 35$
\end{enumerate}

%% ============================================================
\section{Priorite 2 --- Forte valeur ajoutee}
%% ============================================================

\subsection{Texte introductif Section 2 (10 min) \done}

Mock m1 : la Section 2 saute directement de \verb|\section{Sequential...}| a \verb|\subsection{State of the art}|. Ajouter un paragraphe chapeau de 3--4 lignes presentant les trois strategies (permABC-SMC, Over-Sampling, Under-Matching) et le plan de la section. \textit{Fait : paragraphe chapeau ajoute en tete de \texttt{2\_sequential.tex} avec renvois aux sous-sections.}

\subsection{Specifier les hyperparametres du Gaussian toy (10 min) \done}

Mock m2/m9 : donner explicitement $(a, b, s, n, K)$ dans le texte de la Section 3.2 ou en legende de figure. Les valeurs sont dans le code (\texttt{run\_performance\_comparison.py}). \textit{Fait : $a=b=5$, $s=10$, $n=10$, $K=20$ ajoutes apres l'equation dans Section 3.2.}

\subsection{Discuter Figure 2 a la lumiere de $M\to\infty$ (15 min)}

Mock m3 : la figure est la mais le texte ne pointe pas explicitement le comportement prior-like de $\beta$ pour grand $M$ ni la concentration des locaux. Ajouter 2--3 phrases dans le paragraphe qui entoure la figure.

\subsection{Comparaison equilibree ABC-Gibbs (discussion)}

Mock m4 : la comparaison actuelle utilise un modele adversarial. Ajouter une phrase reconnaissant que ABC-Gibbs marche bien quand les dependances sont faibles, ou ajouter un petit benchmark supplementaire.

\subsection{Disambiguer la notation $M$ (10 min) \done}

Mock m6 : Section 2.1 utilise $M$ pour le nombre de datasets par parametre (Bornn et al.), Section 2.3 utilise $M$ pour le nombre de compartiments simules. Renommer l'un des deux (par ex. $M_{\text{sim}}$ dans 2.1 ou garder $M$ seulement pour OS). \textit{Fait : renomme en $M_{\mathrm{rep}}$ dans 2.1 (3 occurrences \texttt{\textbackslash old\{\}} + 1 \texttt{\textbackslash new\{\}}). $M$ reserve a l'Over-Sampling.}

\subsection{Unique particle rate : utiliser ou supprimer (10 min)}

Mock m7 : cette metrique est introduite dans la Section 2.1 mais jamais reportee dans aucune experience. Deux options :
\begin{itemize}[nosep]
    \item La reporter dans au moins une figure (elle est deja calculee dans le code)
    \item Supprimer la discussion (2 phrases)
\end{itemize}

\subsection{Benchmark cout computationnel assignment (code + cluster)}

Mock m8 : pour $K=94$ (SIR), quelle fraction du runtime est dans l'assignment vs. la simulation ? Et quantifier le gain warm-swap vs. LSA complet. Cela renforce la Section 2.3 sur les swaps warm-started.

\subsection{Discussion LFI elargie (30 min)}

Mock M8 + Referee 2 : le paysage LFI actuel inclut SNPE/SNLE (normalizing flows), SBI, methodes bayesiennes generalisees. Le papier ne les mentionne pas. Ajouter un paragraphe dans la conclusion ou l'introduction positionnant permABC :
\begin{itemize}[nosep]
    \item permABC est complementaire : il exploite la structure echangeable, ce que les methodes NN ne font pas nativement
    \item Les methodes amortised (SNPE) sont plus adaptees quand on veut faire de l'inference repetee sur des donnees de meme structure
    \item permABC est plus simple a implementer et ne necessite pas d'entrainement
\end{itemize}

\subsection{Mettre a jour abstract et conclusion (20 min)}

L'abstract et la conclusion ne mentionnent pas :
\begin{itemize}[nosep]
    \item SW$_2$ comme metrique (au lieu de $\varepsilon$)
    \item Le lien formel avec WABC
    \item Les resultats asymptotiques ($K\to\infty$, $M\to\infty$)
    \item La validation SIR synthetique
    \item Les warm-started swaps ($\mathcal{O}(K^2)$ amorti)
\end{itemize}

%% ============================================================
\section{Priorite 3 --- Preparation soumission}
%% ============================================================

\subsection{Cover letter}

JASA exige une cover letter detaillant la significance du travail. Points a couvrir :
\begin{itemize}[nosep]
    \item Motivation : ABC pour modeles hierarchiques echangeables
    \item Contribution : permABC + OS/UM + theorie + SIR 94 departements
    \item Section ciblee : Theory and Methods
    \item Le papier a ete soumis puis rejete par Biometrika ; mentionner que les commentaires des reviewers ont ete integres (sans entrer dans les details)
    \item Pas de conflit d'interet
\end{itemize}

\subsection{Version anonyme}

JASA est double-anonymous. Il faut :
\begin{itemize}[nosep]
    \item Mettre \verb|\anon| a 0 dans \texttt{resubmission.tex}
    \item Verifier qu'aucun nom/affiliation/funding n'apparait dans le PDF
    \item Verifier le supplement aussi
    \item Retirer ou anonymiser les references a ``notre package Python'' et l'URL GitHub
    \item Retirer les commentaires \verb|\al{}|, \verb|\rr{}|, \verb|\xian{}|
\end{itemize}

\subsection{Author Contributions Checklist (ACC)}

Depuis septembre 2021, JASA exige un formulaire ACC pour T\&M. A remplir lors de la revision (pas necessaire pour la soumission initiale, mais preparer le terrain) :
\begin{itemize}[nosep]
    \item Code disponible : oui (GitHub)
    \item Donnees disponibles : oui (Sante publique France, open data)
    \item Instructions de reproduction : oui (\texttt{README.md})
\end{itemize}

\subsection{Reproducibility materials}

Preparer un zip \texttt{reproducibility\_materials.zip} contenant :
\begin{itemize}[nosep]
    \item Le package \texttt{permabc/} + \texttt{requirements.txt} + \texttt{setup.py}
    \item Les scripts de figures (\texttt{experiments/scripts/figures/})
    \item Les donnees COVID (\texttt{experiments/data/})
    \item Un \texttt{README} expliquant comment reproduire chaque figure
\end{itemize}

\subsection{Verifications finales}

\begin{itemize}[nosep]
    \item Abstract $\leq$ 200 mots, pas de references dans l'abstract
    \item Keywords (3--6) pas dans le titre : OK actuellement
    \item Captions sans reference aux couleurs (``dashed'' au lieu de ``blue'')
    \item Toutes les figures dans le corps du texte (pas a la fin)
    \item Pages numerotees
    \item Pas de footnotes avec des maths
\end{itemize}

%% ============================================================
\section{Ordre d'execution recommande}
%% ============================================================

\subsection*{Etape A --- Immediate (cette semaine)}
\begin{enumerate}[nosep]
    \item Lancer le job SIR synthetique sur le cluster (\S\ref{sec:sir_synth})
    \item Envoyer le PDF a Robin et Xian pour validation des \verb|\old|/\verb|\new|
    \item Hardcoder les references supplement (\S\ref{sec:refs})
    \item Redaction rapide : intro Section 2, hyperparametres Gaussian, notation $M$, unique particle rate, abstract/conclusion
\end{enumerate}

\subsection*{Etape B --- Theorie (1--2 semaines)}
\begin{enumerate}[nosep]
    \item Attenuer Theorem 1 + etendre aux ties (\S\ref{sec:thm1}) --- avec Robin/Xian
    \item Formaliser Propositions $M\to\infty$ (\S\ref{sec:minfty})
\end{enumerate}

\subsection*{Etape C --- Numerique (en parallele de B)}
\begin{enumerate}[nosep]
    \item Recuperer resultats SIR synthetique, generer la figure
    \item Creer et lancer le benchmark WABC/Hilbert/Sinkhorn (\S\ref{sec:wabc})
    \item Benchmark cout assignment $K=94$
\end{enumerate}

\subsection*{Etape D --- Finalisation}
\begin{enumerate}[nosep]
    \item Integrer les retours de Robin/Xian, resoudre \verb|\old|/\verb|\new| (\S\ref{sec:oldnew})
    \item Discussion LFI elargie + ABC-Gibbs equilibre
    \item Cover letter, anonymisation, verifications finales
    \item Soumission
\end{enumerate}

\end{document}
